% Compatibility is quarantined from the canonical tutorial and reference.
\section{Compatibility spellings}\label{ch:compatibility}

Compatibility aliases exist to allow older documents to compile without
surprise.  In new manuals, benchmarks, and research papers, always prefer the
canonical modern spellings.  The table below, generated automatically from the
executable language registry, provides the exhaustive alias list.  Since the S4
surface swap, the front ends that parsed legacy dialects have been retired;
the command \tncmd{tenkzkernel} remains available as an inert operation that
rebinds the surface \pkg{} already established at load, so documents that still
include it compile without incident.

\IfFileExists{chapters2/generated-language-aliases.tex}{%
  \noindent% Generated by scripts/tenkz_language.py; do not edit.
{\footnotesize
\begin{tabularx}{\linewidth}{@{}l l X@{}}\toprule Scope & Alias & Replacement \\ \midrule
\bottomrule\end{tabularx}
\medskip
\begin{tabularx}{\linewidth}{@{}l l X@{}}\toprule Scope & Retired spelling & Migration \\ \midrule
kernel-\allowbreak{}mark & \texttt{form=brace-\allowbreak{}above} & use \allowbreak{}form=bracket \allowbreak{}with \allowbreak{}label \allowbreak{}pos=90; \allowbreak{}a \allowbreak{}bracket \allowbreak{}speaks \allowbreak{}from \allowbreak{}the \allowbreak{}side \allowbreak{}its \allowbreak{}label \allowbreak{}sits \allowbreak{}on \\
kernel-\allowbreak{}mark & \texttt{form=brace-\allowbreak{}below} & use \allowbreak{}form=bracket,\allowbreak{} \allowbreak{}which \allowbreak{}speaks \allowbreak{}from \allowbreak{}the \allowbreak{}south \allowbreak{}when \allowbreak{}it \allowbreak{}names \allowbreak{}no \allowbreak{}side \\
kernel-\allowbreak{}mark & \texttt{form=cut} & no \allowbreak{}successor: \allowbreak{}the \allowbreak{}sources' \allowbreak{}cut \allowbreak{}is \allowbreak{}the \allowbreak{}contour \allowbreak{}of \allowbreak{}what \allowbreak{}it \allowbreak{}separates,\allowbreak{} \allowbreak{}so \allowbreak{}use \allowbreak{}form=enclosure \allowbreak{}over \allowbreak{}that \allowbreak{}selection \\
kernel-\allowbreak{}mark & \texttt{form=band} & use \allowbreak{}form=enclosure \allowbreak{}with \allowbreak{}tint,\allowbreak{} \allowbreak{}which \allowbreak{}lays \allowbreak{}ink \allowbreak{}over \allowbreak{}the \allowbreak{}paper \allowbreak{}its \allowbreak{}contour \allowbreak{}encloses \\
command & \texttt{\textbackslash{}tnput} & use \allowbreak{}\textbackslash{}tn[at=(\allowbreak{}r,\allowbreak{}c)],\allowbreak{} \allowbreak{}which \allowbreak{}places \allowbreak{}the \allowbreak{}atom \allowbreak{}at \allowbreak{}the \allowbreak{}address \allowbreak{}the \allowbreak{}put \allowbreak{}named \\
command & \texttt{\textbackslash{}tnsite} & use \allowbreak{}\textbackslash{}tn[at=(\allowbreak{}r,\allowbreak{}c)] \\
command & \texttt{\textbackslash{}tnghost} & no \allowbreak{}successor: \allowbreak{}an \allowbreak{}address \allowbreak{}references \allowbreak{}an \allowbreak{}empty \allowbreak{}cell \allowbreak{}directly \\
command & \texttt{\textbackslash{}tnjoin} & use \allowbreak{}\textbackslash{}tnwire \\
command & \texttt{\textbackslash{}tnedge} & use \allowbreak{}\textbackslash{}tnwire \\
command & \texttt{\textbackslash{}tnarrow} & use \allowbreak{}\textbackslash{}tnwire[dir=to]; \allowbreak{}a \allowbreak{}directed \allowbreak{}index \allowbreak{}is \allowbreak{}the \allowbreak{}wire's \allowbreak{}own \allowbreak{}direction \\
command & \texttt{\textbackslash{}tncut} & use \allowbreak{}\textbackslash{}tnmark[form=enclosure] \allowbreak{}over \allowbreak{}the \allowbreak{}selection \allowbreak{}a \allowbreak{}cut \allowbreak{}separates; \allowbreak{}the \allowbreak{}cut \allowbreak{}form \allowbreak{}is \allowbreak{}retired \allowbreak{}with \allowbreak{}it \\
command & \texttt{\textbackslash{}tnregion} & use \allowbreak{}\textbackslash{}tnmark[form=enclosure] \\
\bottomrule\end{tabularx}
}
%
}{%
  \begin{quote}\small
  The generated alias table is not present.  Generate it from the executable
  language registry before publishing the manual.
  \end{quote}
}

Whenever an alias is parsed, it normalizes to the exact same semantic events as
its canonical replacement, without altering value types, defaults, scopes, or
diagnostic codes.  The canonical-only linter is authoritative for all benchmark
sources: when updating legacy code, replace the alias with its canonical modern
counterpart.  Historical notes in \tnfile{docs/tenkz/history/} serve as archival
records and do not define supported aliases.

Notice that there is no compatibility form of \tncmd{tndefine}: it was once
proposed but never implemented in the language, so no alias is needed.  For
custom typed atoms, use \tncmd{tndeclareatom}.
